% from ega3-7-fr.tex
\oldpage[III]{182}

\begin{corollary}[7.3.4]
Supposons $A$ commutatif, $T_p$ exact et supposons en outre que $T_p(A)$ soit
un $A$-module de présentation finie. Alors la fonction
\[
  x\mapsto\operatorname{rang}_{\kappa(x)}(T_p(\kappa(x)))
\]
est localement constante dans $X=\operatorname{Spec}(A)$, donc constante si
$\operatorname{Spec}(A)$ est connexe.
\end{corollary}

\begin{proof}
En effet, comme $T_p(A)$ est un $A$-module plat en vertu de (7.3.3, $f)$), il
est projectif de type fini, et $\bigl(T_p(A)\bigr)^{\sim}$ est donc un
$\mathscr O_X$-Module localement libre (Bourbaki, \emph{Alg. comm.}, chap. II,
\S 5, n\textsuperscript{o} 2, th. 1); on a en outre
\[
  T_p(\kappa(x))=T_p(A)\otimes_A\kappa(x)
\]
(7.3.3, $e)$) et l'on sait que le rang au point $x$ du $A$-module $T_p(A)$ est
localement constant (\emph{loc. cit.}), d'où le corollaire.
\end{proof}

\begin{proposition}[7.3.5]
Supposons que $A$ soit un anneau artinien à gauche dont le quotient par son
radical $\mathfrak m$ est un corps $k$. Alors les conditions de (7.3.3) sont
encore équivalentes à chacune des suivantes :
\begin{enumerate}
  \item[g)] L'homomorphisme canonique
    \[
      T_i(\varepsilon):T_i(A_s)\longrightarrow T_i(k)
    \]
    est un épimorphisme pour $i=p$ et $i=p+1$.
  \item[h)] $T_p(\varepsilon)$ est un épimorphisme et $T_p(A_s)$ est un
    $A$-module à droite plat (ou, ce qui revient au même (Bourbaki,
    \emph{Alg. comm.}, chap. II, \S 3, n\textsuperscript{o} 2, cor. 2 de la
    prop. 5) un $A$-module libre).
\end{enumerate}
Supposons en outre que $A$ soit commutatif et le $A$-module $T_p(k)$ de
longueur finie $d$. Alors les conditions précédentes équivalent aussi à
chacune des suivantes :
\begin{enumerate}
  \item[i)] Pour tout $A$-module $M$ de longueur finie, on a
    \[
      \operatorname{long}(T_p(M))=d\mathbin{.}\operatorname{long}(M).
      \tag{7.3.5.1}\label{III.7.3.5.1.fr}
    \]
  \item[j)] On a
    \[
      \operatorname{long}(T_p(A))=d\mathbin{.}\operatorname{long}(A).
      \tag{7.3.5.2}\label{III.7.3.5.2.fr}
    \]
\end{enumerate}
\end{proposition}

L'équivalence de $g)$ et $h)$ avec les conditions de (7.3.3) découle aussitôt
de (7.2.7). Pour démontrer les autres assertions, nous utiliserons le lemme
suivant :

\begin{lemma}[7.3.5.3]
Soient $\mathcal K$, $\mathcal K'$ deux catégories abéliennes,
$F:\mathcal K\to\mathcal K'$ un foncteur
covariant additif; on suppose que $F$ soit semi-exact, et que, pour tout objet
simple $S$ de $\mathcal K$, $F(S)$ soit un objet de longueur finie dans
$\mathcal K'$. Alors, pour tout objet $E$ de longueur finie dans $\mathcal K$,
$F(E)$ est de longueur finie dans $\mathcal K'$. Pour toute suite exacte
\[
  0\longrightarrow E'\xrightarrow{u}E\xrightarrow{v}E''\longrightarrow0
\]
d'objets de longueur finie dans $\mathcal K$, on a
\[
  \operatorname{long}F(E)
  \leq
  \operatorname{long}F(E')+\operatorname{long}F(E'')
  \tag{7.3.5.4}\label{III.7.3.5.4.fr}
\]
et pour que les deux membres de (7.3.5.4) soient égaux, il faut et il suffit
que la suite
\[
  0\longrightarrow F(E')\longrightarrow F(E)\longrightarrow F(E'')
  \longrightarrow0
\]
soit exacte.
\end{lemma}

\begin{proof}
En effet, la suite
\[
  F(E')\xrightarrow{F(u)}F(E)\xrightarrow{F(v)}F(E'')
\]
est exacte par hypothèse; si l'on suppose $F(E')$ et $F(E'')$ de longueur
finie, il en est de même de $\operatorname{Im}(F(u))$ et de
$\operatorname{Im}(F(v))$ et comme
$\operatorname{Ker}(F(v))=\operatorname{Im}(F(u))$, $F(E)$ est de longueur
finie et l'on a
\[
  \operatorname{long}F(E)
  =
  \operatorname{long}\operatorname{Im}(F(u))
  +
  \operatorname{long}\operatorname{Im}(F(v))
  \leq
  \operatorname{long}F(E')+\operatorname{long}F(E'').
  \tag{7.3.5.5}\label{III.7.3.5.5.fr}
\]
